\documentclass{article}
\usepackage{amsmath,amssymb}
\usepackage[vocab={realfunctions,asymptotics}]{leanify}
\newtheorem{lemma}{Lemma}
\begin{document}

\section{Notation}

We write $L(f)$ for the limit of $f$ at infinity.
We write $\rho(f)$ for the index of regular variation of $f$.
For every real number $x$, let $q(x)\coloneqq x^2+1$.

\section{Comparisons}

\begin{lemma}\label{equivalent}
Let $f$ and $g$ be functions. If $f$ is asymptotically equivalent to $g$, then $f$ is dominated by $g$.
\end{lemma}

\begin{lemma}\label{negligible}
Let $f$, $g$ and $h$ be functions. If $f$ is negligible compared to $g$ and $g$ is dominated by $h$, then $f$ is negligible compared to $h$.
\end{lemma}

\begin{lemma}\label{order}
Let $f$ and $g$ be functions. If $f$ is of the same order as $g$, then $g$ is of the same order as $f$.
\end{lemma}

\begin{lemma}\label{bounded}
Let $f$ be a function bounded at infinity. Then $f$ is of polynomial growth.
\end{lemma}

\begin{lemma}\label{dominated}
Let $f$ and $g$ be functions. If $f$ is dominated by $g$ and $g$ is bounded at infinity, then $f$ is bounded at infinity.
\end{lemma}

\begin{lemma}\label{polynomial}
$q$ is of polynomial growth and $q$ is of at most exponential growth.
\end{lemma}

\section{Regular variation and limits}

\begin{lemma}\label{slow}
Every slowly varying function is eventually positive.
\end{lemma}

\begin{lemma}\label{regular}
Let $\alpha$ be a real number and let $f$ be a function regularly varying with index $\alpha$. Then $\rho(f)=\alpha$.
\end{lemma}

\begin{lemma}\label{limit}
Let $f$ and $g$ be functions. Suppose that $f$ is asymptotically equivalent to $g$ and that $L(g)=1$. Then $L(f)=1$.
\end{lemma}

\end{document}
